% ─── Описание моделей ─────────────────────────────────────────────────────
\begin{frame}[shrink=10]{Исследуемые модели}

  \begin{columns}[T]
    \column{0.32\textwidth}
    \begin{tcolorbox}[infobox, equal height group=models]
      \textbf{\color{primary}\faChartLine\ ARIMA}\\[6pt]
      \small
      \textcolor{accent}{\faAngleRight}~Классическая статистическая модель\\[3pt]
      \textcolor{accent}{\faAngleRight}~Авторегрессия + скользящее среднее\\[3pt]
      \textcolor{accent}{\faAngleRight}~Дифференциация для стационарности\\[3pt]
      \textcolor{accent}{\faAngleRight}~Два режима: rolling и static
    \end{tcolorbox}

    \column{0.32\textwidth}
    \begin{tcolorbox}[infobox, equal height group=models]
      \textbf{\color{primary}\faTree\ Random Forest}\\[6pt]
      \small
      \textcolor{accent}{\faAngleRight}~Ансамбль деревьев решений\\[3pt]
      \textcolor{accent}{\faAngleRight}~Bagging + случайные признаки\\[3pt]
      \textcolor{accent}{\faAngleRight}~Устойчивость к шуму\\[3pt]
      \textcolor{accent}{\faAngleRight}~Лаговые признаки на входе
    \end{tcolorbox}

    \column{0.32\textwidth}
    \begin{tcolorbox}[infobox, equal height group=models]
      \textbf{\color{primary}\faBrain\ LSTM}\\[6pt]
      \small
      \textcolor{accent}{\faAngleRight}~Рекуррентная нейросеть\\[3pt]
      \textcolor{accent}{\faAngleRight}~Шлюзы: forget, input, output\\[3pt]
      \textcolor{accent}{\faAngleRight}~Долгосрочная память\\[3pt]
      \textcolor{accent}{\faAngleRight}~50 нейронов, окно = 10 дней
    \end{tcolorbox}
  \end{columns}

  \vspace{0.3cm}
  \begin{tcolorbox}[successbox]
    \small \textbf{Ключевой вопрос:} способны ли нейросети с нелинейным моделированием превзойти проверенные временем статистические методы на реальных данных курса USD/UZS?
  \end{tcolorbox}
\end{frame}
